\section{Conclusão}

\hspace{0.4cm}O presente estudo demonstrou que a ferramenta computacional desenvolvida é eficaz para análises comparativas com as aplicações numéricas propostas, produzindo resultados satisfatórios e totalmente coerentes para a avaliação de perdas de protensão em estruturas de concreto. Além de seu desempenho técnico, destaca-se sua fácil acessibilidade e a capacidade de esclarecer os cálculos envolvidos, graças ao repositório aberto, à presente documentação e à presença de exemplos bem elaborados e comparativos. A interface foi planejada com clareza e simplicidade, exibindo informações de forma organizada para facilitar o uso por usuários sem conhecimentos em programação.

Espera-se, portanto, que a ferramenta contribua para uma maior capacidade analítica dos processos de cálculo nos projetos de estruturas de concreto protendido, incentivando otimizações mais precisas e seguras. Além disso, é fundamental ressaltar novamente a importância das análises de perdas de protensão, que são cruciais para garantir a segurança e a eficiência das estruturas de concreto protendido. Em relação a estudos e projetos futuros, destaca-se a possibilidade de evolução da própria ferramenta, que, por ser de código aberto, está aberta a contribuições e melhorias contínuas, tanto da comunidade de usuários quanto de outros desenvolvedores interessados.